%% KEY ADDITIONS for paper restructure
%% Each block goes into the specified section

%% ============================================================
%% FOR SECTION 2: D(1/p) Proposition
%% ============================================================

\begin{proposition}[Dominance of $1/p$]\label{prop:D1p}
For every prime $p\ge3$, the rank deviation of the fraction $1/p$ satisfies
\[
  D\!\left(\tfrac{1}{p}\right) = 1 - \frac{|{\F}_{p-1}|}{p}
  \;=\; 1 - \frac{1+\sum_{k=1}^{p-1}\varphi(k)}{p}
  \;\sim\; -\frac{3p}{\pi^2}
  \quad(p\to\infty).
\]
Moreover, the product $D(1/p)\cdot\delta(1/p) = D(1/p)/p \to -3/\pi^2\approx -0.304$
is a universal constant, independent of~$p$.  Computationally, this single fraction
accounts for $65{-}73\%$ of the total $|\sum D\cdot\delta|$ for $p\le199$.
\end{proposition}

\begin{proof}
Since $p\ge3$ is prime, $1/(p{-}1) > 1/p$, so $1/p$ lies strictly between
$0/1$ and $1/(p{-}1)$ in ${\F}_{p-1}$.  The only fraction in ${\F}_{p-1}$
at most $1/p$ is $0/1$, giving $\rank(1/p, {\F}_{p-1}) = 1$.
The formula follows.  The asymptotic uses
$\sum_{k=1}^{N}\varphi(k) = 3N^2/\pi^2 + O(N\log N)$ (Mertens, 1874).
\end{proof}

%% ============================================================
%% FOR SECTION 3: Damage/Response Decomposition
%% ============================================================

\begin{observation}[Damage--Response decomposition]\label{obs:damage-response}
Define the \emph{damage ratio} $R_1(p)$ by summing $D\cdot\delta_1$
over new fractions $a/p$ entering ${\F}_p$, and the \emph{response ratio}
$R_2(p)$ by summing $D\cdot\delta_2$ over existing fractions $a/b$ in
${\F}_{p-1}$, where $\delta_2(a/b) = (a - pa\bmod b)/b$ is the
insertion deviation.

For all $4{,}617$ primes with $M(p)\le-3$ and $p\le100{,}000$:
\begin{itemize}
\item $R_1(p) < 0$ always (new fractions \emph{damage} regularity);
\item $R_2(p) > 0$ always (existing fractions shift to
  \emph{overcompensate} the damage).
\end{itemize}
The net effect $\Delta W(p) < 0$ (regularity improves) arises because
the response dominates the damage.
\end{observation}

\begin{remark}
The sign alignment underlying $R_2>0$ is weak:
for $p=31$, only $59\%$ of old fractions have
$\sgn(D)\cdot\sgn(\delta_2)>0$.  The positivity of
$\sum D\cdot\delta_2$ is driven primarily by magnitude:
fractions with large positive $D$ also have large positive $\delta_2$,
contributing disproportionately.  An analytical proof of $R_2>0$
remains open.
\end{remark}

%% ============================================================
%% FOR SECTION 3 or 6: Composites Characterization
%% ============================================================

\begin{observation}[Healing characterization]\label{obs:composites}
For $n\le 500$:
\begin{enumerate}
\item Every composite $n$ with $\omega(n)\ge 2$
  (two or more distinct prime factors) satisfies $\Delta W(n) > 0$.
\item The non-healing composites ($\Delta W(n) \le 0$) are
  precisely the prime powers $p^k$.
\end{enumerate}
The number-theoretic boundary between damage and healing is exactly
$\omega(n) = 1$ versus $\omega(n)\ge 2$.
\end{observation}

%% ============================================================
%% FOR SECTION 6: GK Concentration
%% ============================================================

\begin{observation}[Gauss--Kuzmin concentration]\label{obs:gk-conc}
For each qualifying prime~$p$, sort the new fractions $a/p$
by their terminal continued-fraction quotient $c_m(a/p)$
in decreasing order.  The top $20\%$ of fractions
(those with the largest $c_m$) contribute over $93\%$
of the total $|\sum D\cdot\delta|$.  This heavy-tail concentration
is stable across all tested primes ($p=31$ to $199$) and reflects
the infinite mean of the Gauss--Kuzmin distribution.
\end{observation}

%% ============================================================
%% FOR ABSTRACT: Updated last paragraph
%% ============================================================

% ADD after current abstract content:
% We introduce the \emph{Farey spectroscope}, a spectral function
% $F(\gamma) = |\sum_p R_2(p)\,p^{-1/2-i\gamma}|^2$ whose peaks
% detect the locations of nontrivial zeta zeros from Farey data alone.
% The first zero $\gamma_1 = 14.135$ is recovered at $14.05$ ($0.6\%$ error)
% from $2{,}729$ qualifying primes.  Twisting by Dirichlet characters
% extends detection to $L$-function zeros.  Predicted and observed
% peak amplitudes correlate at $r = 0.953$.
